\section{Implementation}
\label{sec:impl}

%% tab-impl.tex -- Table (single column): the implementation, by the numbers.
%%
%% Every figure in this table comes from paper/generated/stats.tex, which
%% paper/scripts/gen_stats.py recomputes from the source tree. The macros are
%% never typed by hand, so the table cannot drift from the artifact.

\begin{table}[tbp]
\centering
\caption{The implementation, by the numbers. Every count here is regenerated from
the source tree by \code{paper/scripts/gen\_stats.py} at build time rather than
transcribed, so it can be recomputed instead of trusted; the tree state is
recorded as \code{\statGitRev}. Line counts are physical lines including comments
and blanks, which is the cheapest reproducible definition rather than the most
flattering one.}
\label{tab:impl}
\small
\renewcommand{\arraystretch}{1.24}
\begin{tabular}{@{}L{8.50cm}R{2.30cm}R{2.30cm}@{}}
\toprule
\theadrow
\thead{Component} & \thead{Files} & \thead{Lines} \\
\midrule
\zebra Python package (control plane)  & \statPyFiles & \statPyLoc \\
TypeScript / TSX (SPA)             & \statUiFiles    & \statUiLoc \\
\zebra Test modules & \statTestFiles & \statTestLoc \\
\midrule
\tgroup{3}{Interface and schema surface} \\
Route modules / HTTP declarations   & \statRouteModules & \statRouteDecls \\
\zebra ORM models / request-response schemas & \statModels & \statSchemas \\
Alembic revisions / workflow node types & \statMigrations & \statNodeTypes \\
\addlinespace[0.2ex]
\tgroup{3}{Architectural constants} \\
Layers / lifecycle modes            & \statLayers     & \statModes \\
\zebra Asset families / RBAC scopes & \statFamilies & \statScopes \\
Background loops / optimizer paths  & \statLoops      & \statOptimizers \\
\bottomrule
\end{tabular}
\end{table}


\subsection{Scale and composition}

\tabref{tab:impl} gives the implementation by the numbers. Every count in that
table --- and every count in this paper's prose --- is produced by a script that
walks the source tree at build time (\code{paper/scripts/gen\_stats.py}) and emits
LaTeX macro definitions that the document consumes. Nothing is transcribed by
hand, so a number cannot drift out of agreement with the artifact it describes, and
a reader with the repository can recompute all of them. Line counts are physical
lines including comments and blanks: the cheapest reproducible definition rather
than the most flattering one.

The backend is Python over Starlette and SQLAlchemy with Alembic migrations; the
frontend is a React single-page application built separately with Vite and served
\emph{through} the Python package from the same origin, so both deployment
topologies present one origin to the browser (D10, \tabref{tab:decisions}).

\subsection{Three implementation choices worth recording}

\para{Synchronous ORM behind async routes.} Route callables are predominantly
async while the ORM is synchronous; selected blocking work is explicitly offloaded
to a thread. The reasoning is D6 in \tabref{tab:decisions}: the workload is
transaction-bound rather than connection-bound, and the synchronous ORM is the more
mature surface. The cost is a standing review obligation --- every new blocking call
site is a latent event-loop stall --- and we record it as a cost rather than
presenting the choice as free.

\para{A deterministic fake behind every provider seam.} \code{LLMProvider},
\code{EvalProvider}, and \code{Promoter} each have a deterministic fake as the
default, so the server boots and the full control plane is exercisable with no API
key configured. Two consequences follow that are larger than the decision looks.
The test suite is hermetic: no test requires a model provider, so
\statTestFiles{} test modules run without network access or spend. And performance
measurement is separable: control-plane latency can be measured independently of
provider latency, which is what makes the protocol in \appref{app:protocol}
meaningful rather than a measurement of somebody else's inference service.

\para{Migrations as an ownership boundary.} \statMigrations{} Alembic revisions
manage \sys's schema, and \sys's database URL is configured independently of
\mlflow{}'s backend store. The bundled development stack deliberately uses separate
logical databases so that the two Alembic histories never compete for one version
table (\figref{fig:topologies}). This is a small operational decision that prevents
a specific and hard-to-diagnose failure --- two migration tools disagreeing about
one schema-version row.

\subsection{Testing and quality gates}

The suite spans three layers: backend tests under pytest with parallel execution
and a separate integration marker for tests requiring external servers; frontend
unit tests; and browser end-to-end tests. Static analysis runs a linter and a type
checker over the package. All three layers can emit a common structured report
format, and emitting results requires no JVM even though rendering the report does.
The browser layer includes a separate Cookbook configuration that starts a migrated,
isolated database and rejects adapters containing direct request, endpoint, database,
or manifest-seeding shortcuts.  Its executable journeys cover skill authoring and
trigger tests, all-sixteen catalog visibility plus installation of every example as
a separate paused draft, and review-queue creation.  The local CI mirror terminates the detached
MLflow/uvicorn process group before deleting its state so a passing browser run does
not leave a listener or database behind.

We are deliberately not reporting a test count or a coverage percentage as evidence
of quality. Both are available, and neither supports the claims this paper makes;
presenting them as though they did would be the kind of number-decoration
\secref{sec:eval} declines to engage in. What the suite does support is the
hermeticity claim above: the control plane is exercisable end to end without a model
provider, which is a checkable structural property rather than a quality score.

\subsection{Deployment}

\figref{fig:topologies} gives both topologies. A bundled loopback-only container
stack provides PostgreSQL~17 with pgvector and Apache~AGE, S3-compatible object
storage, \mlflow{}, an AI gateway, and NATS, for development. It is explicitly not
a production deployment template: it is loopback-bound, and it enables a
convenience credential bootstrap that a network-reachable deployment must leave
disabled. Saying so in the paper is not throat-clearing --- a development stack
mistaken for a production template is a realistic and severe failure mode.
